\def\probtitle{XOR 팰린드롬}
\def\probno{K} % 문제 번호

\begin{problem}{\probtitle}{표준 입력(stdin)}{표준 출력(stdout)}{2 초}{1024 megabytes}

길이가 $M$인 이진 문자열 $N$개가 주어진다.

주어진 문자열 중 1개 이상을 선택하여, 선택한 모든 문자열을 비트별 배타적 논리합(bitwise XOR)한 결과 문자열이 팰린드롬이 되는 경우의 수를 구하여라.

단, 문자열을 선택하는 순서는 고려하지 않는다. 즉, 고르는 순서만 다르고 선택된 문자열의 구성이 같은 경우는 한 번만 센다. 또한, 입력으로 같은 이진 문자열이 여러 번 주어지더라도, 이들은 서로 구별되는 다른 문자열로 취급한다.

배타적 논리합(bitwise XOR) 또는 팰린드롬에 대한 설명은 하단의 노트를 참고하라.

\InputFile

첫째 줄에 정수 $N$, $M$이 공백으로 구분되어 주어진다.

둘째 줄부터 $N$개의 줄에 걸쳐 길이가 $M$인 이진 문자열 $T_1, T_2, \dots, T_N$이 한 줄에 하나씩 주어진다.

\OutputFile

조건을 만족하도록 문자열을 선택하는 경우의 수를 출력한다. 단, 답이 매우 커질 수 있으므로 $1\, 000\, 000\, 007(=10^9+7)$로 나눈 나머지를 출력한다.

\Constraints

\begin{itemize}[topsep=0pt,noitemsep]
    \item $1 \le N \le 50\, 000$
    \item $1 \le M \le 20$
    \item $1 \le i \le N$ 인 모든 $i$에 대해 $|T_i| = M$ 이다.
    \item $1 \le i \le N$ 인 모든 $i$에 대해 $T_i$는 \texttt{0}과 \texttt{1}로만 이루어져 있다.
\end{itemize}

\Example

\begin{example}
    \exmpfile{./example/01.in.txt}{./example/01.out.txt}%
    \exmpfile{./example/02.in.txt}{./example/02.out.txt}%
\end{example}

\Notes

두 이진 문자열 $A, B$의 비트별 배타적 논리합(bitwise XOR) $A \oplus B$는 다음과 같이 정의된다.

\begin{itemize}[topsep=0pt,noitemsep] 
\item $A$의 $i$번째 글자와 $B$의 $i$번째 글자가 서로 다르면 $A \oplus B$의 $i$번째 글자가 $1$이고, 같으면 $0$이다.
\end{itemize}

예를 들어 $\text{1100} \oplus \text{1010}$은 $\text{0110}$으로 계산한다.

$k$개의 이진 문자열 $S_1, S_2, \cdots, S_k$를 비트별 배타적 논리합한 결과는 $(\dots((S_1 \oplus S_2) \oplus S_3) \oplus \dots) \oplus S_k$로 정의된다. 이 연산의 결과는 $S_1, S_2, \cdots, S_k$의 순서를 바꾸더라도 변하지 않음을 증명할 수 있다.

팰린드롬이란, 앞에서부터 읽으나 뒤에서부터 읽으나 같은 문자열을 말한다. 예를 들어, \texttt{010}, \texttt{0110} 은 팰린드롬이며 \texttt{011}, \texttt{1011} 은 팰린드롬이 아니다.

\end{problem}